\documentclass[a4paper,11pt]{article}
\usepackage[utf8]{inputenc}
\usepackage[T1]{fontenc}
\usepackage[french]{babel}
\usepackage[left=2cm,right=2cm,top=2cm,bottom=2cm]{geometry}
\usepackage{xcolor}
\usepackage{hyperref}
\usepackage{fontawesome5}
\usepackage{parskip}

\definecolor{primary}{RGB}{200, 50, 30}
\hypersetup{colorlinks=true, urlcolor=primary}

\begin{document}

\noindent
\begin{minipage}[t]{0.5\textwidth}
    \textbf{Loïc Aron MBASSI EWOLO}\\
    Élève-Ingénieur Bac+5 -- Double diplôme ENSEA\\[3pt]
    \faMapMarker\ Cergy, France\\
    \faPhone\ (+33) 7 59 52 33 98\\
    \faEnvelope\ \href{mailto:loicaron.mbassiewolo@ensea.fr}{loicaron.mbassiewolo@ensea.fr}
\end{minipage}
\hfill
\begin{minipage}[t]{0.4\textwidth}
    \raggedleft
    Cergy, le 10 mars 2026
\end{minipage}

\vspace{1.2cm}

\hfill
\begin{minipage}[t]{0.5\textwidth}
    \raggedleft
    \textbf{Framatome}\\
    Service Recrutement\\
    Grenoble, France
\end{minipage}

\vspace{1cm}

\noindent\textbf{Objet :} Candidature en alternance -- Développeur logiciel embarqué

\vspace{0.8cm}

Madame, Monsieur,

Concevoir un schéma de détection et d'isolation de défauts sur drone, puis implémenter un contrôle tolérant aux pannes avec reconfiguration automatique de contrôleur : ce projet à l'ENSEA m'a confronté aux exigences des systèmes safety-critical, où chaque ligne de code doit être spécifiée, testée et documentée. C'est cette rigueur que je souhaite mettre au service du développement de logiciels de sûreté nucléaire chez Framatome.

Mon expérience en développement embarqué C/C++ couvre plusieurs contextes exigeants. Dans le challenge CoHoMa (robotique autonome militaire, ENSEA/Armée française), je développe en C/C++ sur Nvidia Jetson avec intégration capteurs (Lidar 3D, caméra de profondeur) et containerisation Docker pour reproductibilité. Mon projet de simulation multiprocessus en C sous Linux (mémoire partagée, signaux POSIX, ordonnanceur FIFO) m'a donné une maîtrise du bas niveau : gestion manuelle de la mémoire, synchronisation de processus, rigueur de programmation conforme aux contraintes de fiabilité. Ces compétences sont directement transposables au codage en C embarqué et aux revues de code selon les normes MISRA et CERT.

En stage à INRIA Saclay (équipe TRIBE, avril à août 2026), je développe des pipelines Python structurés incluant de l'analyse statique de code et de la documentation technique rigoureuse en anglais. Ce travail complète mon profil par la méthodologie de tests et de validation logicielle que requiert votre environnement qualité. Mon prototypage IoT sur STM32 (capteurs IMU, fusion de données, soudure) au laboratoire IOTA de l'École Polytechnique de Yaoundé témoigne de ma capacité à intervenir au niveau hardware comme software.

Je suis disponible dès septembre 2026 pour une alternance de 12 mois à Grenoble. Découvrir le développement de logiciels dédiés à la sûreté nucléaire, de la rédaction des spécifications à la validation, est une opportunité technique et humaine que je souhaite saisir.

Je reste à votre disposition pour un entretien afin de discuter de ma candidature.

Je vous prie d'agréer, Madame, Monsieur, l'expression de mes salutations distinguées.

\vspace{0.8cm}

\textbf{Loïc Aron MBASSI EWOLO}\\
\textit{Élève-Ingénieur ENSEA / Polytechnique Yaoundé}

\end{document}
